% V7 results draft based on the audited DFG archived score report.
% Reference: OXTS/GPS correlated, NOT independent physical GT.
\section{Experimental Evaluation}
\label{sec:experiments_v7}

\subsection{Development and Reserved Sequence Protocol}
The reliability networks were trained on a synthetic degraded-GPS
candidate extracted from KITTI Raw sequence
\texttt{2011\_10\_03\_drive\_0027\_sync} (4,544 frames).
The training anchors were $[63,3178)$ (3,115 samples)
and the previously explored development-validation anchors were
$[3180,3860)$ (680 samples). The 0027 validation partition
cannot be treated as an independent test.

Cross-date evaluation used
\texttt{2011\_09\_30\_drive\_0016\_sync} (279 frames) and
\texttt{2011\_09\_30\_drive\_0020\_sync} (1,104 frames).
Synthetic degradation is generated from a fixed development
schedule resampled in normalized sequence time. In each drive,
all four compared methods use exactly the same degraded GPS,
gyroscope, LiDAR and stereo factor measurements and calibration.
Only GPS position-factor covariance differs.

\subsection{Evaluation Metrics and Compared Methods}
Let $\mathbf{p}_t$ be optimized pose translation and
$\mathbf{q}_t$ the original \emph{uncorrupted OXTS-derived GPS}
position. The primary origin-relative three-dimensional error is
\begin{equation}
E_{3D}=\sqrt{\frac{1}{N}
\sum_{t=0}^{N-1}
\|\mathbf{p}_t-(\mathbf{q}_t-\mathbf{q}_0)\|_2^2}.
\label{eq:v7_eval_rmse}
\end{equation}
The 2D version uses the first two axes. Translation-only
$k$-frame relative displacement RMSE ($k=1,10$)
compares $\mathbf{p}_{t+k}-\mathbf{p}_t$ with
$\mathbf{q}_{t+k}-\mathbf{q}_t$. These are neither
six-degree-of-freedom RPE nor independent physical
ground-truth localization errors: original OXTS GPS
also generated the synthetic degraded GPS measurements.

The four fixed cases are Axis-Fixed (training-only
axis RMS covariance), GRU96-Cap4, GRU106-Cap4
(capacity control), and Mamba-Cap4. The three
learned cases apply the same training-fitted five-bin
mapping and GPS information cap; only the $h=0$
current-reliability head is used to weight the graph.

\subsection{Cross-Date Position Error}
Table~\ref{tab:v7_main} reports primary position errors.
Mamba-Cap4 has the smallest 3D error on 0016 at
2.4467 m, narrowly below GRU96-Cap4 at 2.4732 m
by 0.0265 m (approximately 1.07\%). On 0020
the ranking reverses: GRU96-Cap4 obtains 3.2032 m
while Mamba-Cap4 obtains 6.9844 m, roughly 2.18
times the GRU96 error.

\begin{table}[!t]
\centering
\caption{Origin-relative OXTS/GPS 3D RMSE (m; lower is better).
The reference is measurement-correlated.}
\label{tab:v7_main}
\begin{tabular}{lcc}
\hline
Method & 0016 (279) & 0020 (1104)\\
\hline
Axis-Fixed & 2.6526 & 21.0158\\
GRU96-Cap4 & 2.4732 & \textbf{3.2032}\\
GRU106-Cap4 & 2.6329 & 4.2117\\
Mamba-Cap4 & \textbf{2.4467} & 6.9844\\
\hline
\end{tabular}
\end{table}

Against Axis-Fixed on 0020, 3D RMSE decreases
by 84.76\% with GRU96-Cap4 and 66.77\% with
Mamba-Cap4. On 0016, the corresponding reductions
are approximately 6.76\% and 7.76\%, respectively.
The improvements indicate a conditional effect
of the tested dynamic GPS weighting. They do not
establish universal Mamba superiority, calibrated
reliability, or a statistically significant difference
between model architectures.

\subsection{Horizontal and Short-Term Motion Diagnostics}
Table~\ref{tab:v7_secondary} compares 2D RMSE and
translation-displacement errors. Mamba-Cap4 yields
the lowest secondary errors on 0016, while
GRU96-Cap4 is best on 0020. On the longer drive,
the 10-frame translation error is 2.0928 m for
GRU96-Cap4, compared with 4.8892 m for Mamba-Cap4
and 9.3911 m for Axis-Fixed.

\begin{table*}[!t]
\centering
\caption{Position and displacement errors (m; lower is better).
Displacement metrics compare translation only, not full SE(3) RPE.}
\label{tab:v7_secondary}
\begin{tabular}{llccc}
\hline
Drive & Method & 2D RMSE & 1-frame RMSE & 10-frame RMSE\\
\hline
0016 & Axis-Fixed & 2.6154 & 0.2778 & 1.6519\\
0016 & GRU96-Cap4 & 2.4323 & 0.2987 & 1.7829\\
0016 & GRU106-Cap4 & 2.5957 & 0.2995 & 1.7824\\
0016 & Mamba-Cap4 & \textbf{2.4209} & \textbf{0.2719} &
\textbf{1.5997}\\
\hline
0020 & Axis-Fixed & 16.5548 & 1.0455 & 9.3911\\
0020 & GRU96-Cap4 & \textbf{2.8535} & \textbf{0.3435} &
\textbf{2.0928}\\
0020 & GRU106-Cap4 & 3.4844 & 0.4529 & 3.2205\\
0020 & Mamba-Cap4 & 6.5991 & 0.5891 & 4.8892\\
\hline
\end{tabular}
\end{table*}

\subsection{Optimization Stability and Execution Audit}
The 0016 graphs share 279 GPS factors, 278 IMU,
278 LiDAR, 275 camera factors, and one anchor prior.
The 0020 graphs share 1,104 GPS, 1,103 IMU,
1,103 LiDAR, 1,063 camera factors and one prior.
All eight runs satisfy the predeclared numerical
stability criterion. Every 0016 method takes two
LM passes. On 0020, GRU96-Cap4 takes three
passes and the other cases take two. Objective
values cannot be compared across methods directly
because the covariance weights differ; numerical
stability does not certify a global optimum.
The recorded LM times omit network inference and
odometry construction and are not a real-time
performance benchmark.

\subsection{Post-Release Synchronization Incident and Limitations}
The initial one-time reserved release built the
0016 four-sensor measurements, but the 0020
preprocessing was interrupted by a strict
100-ms sensor-to-OXTS timestamp requirement.
A read-only investigation after the release
identified common negative startup offsets
of approximately 123--126 ms at precisely
frames 0--2 of the 0020 LiDAR/stereo streams.
A narrow \emph{post-release} engineering exception
was then introduced for this shared early-frame
pattern. Neither frames nor timestamps were
deleted, reordered or changed; the original
0016 build files were preserved. Both drives
were solved and scored only after the audited
recovery, and the final report and manifests
were subsequently verified with SHA256.

Because the exception was defined after the
reserved drives were opened, these results
cannot be described as fully preregistered or
entirely untouched. Although reference
position errors did not inform the noise rules
or model selection, this procedural limitation
must be disclosed. Additional limitations are
the measurement-correlated OXTS reference,
only two cross-date sequences, one model seed,
one held-out initialization, and the absence
of incremental inference or full-state accuracy
evaluation. In particular, the predictive
$h=1,2$ and Beta evidence uncertainty heads
are not used by the evaluated factor graph.
The experimental evidence therefore concerns
\emph{current-reliability-guided adaptive GPS
covariance}, not validated anticipatory
Mamba-based sensor fusion.
